\chapter{Datasets}
\label{app:datasets}

\section{CICIDS2017}

The Canadian Institute for Cybersecurity Intrusion Detection System
dataset 2017 \cite{sharafaldin2018} contains network traffic captures
over five days with 14 attack classes.
Features are extracted using CICFlowMeter from PCAP captures.

\begin{table}[H]
\centering
\caption{CICIDS2017 class distribution (approximate).}
\begin{tabular}{@{}lrr@{}}
\toprule
Attack Class & Count & Fraction \\
\midrule
BENIGN           & 2,273,097 & 53.0\% \\
DDoS             & 128,027   & 15.5\% \\
PortScan         & 158,930   & 14.0\% \\
DoS Hulk         & 231,073   &  6.0\% \\
DoS GoldenEye    &  10,293   &  2.0\% \\
Bot              &   1,966   &  2.2\% \\
Web Attack       &   2,180   &  1.9\% \\
SSH-Patator      &   5,897   &  1.0\% \\
FTP-Patator      &   7,938   &  0.9\% \\
Heartbleed       &      11   &  0.01\% \\
Others           & --- & --- \\
\bottomrule
\end{tabular}
\end{table}

\textbf{Download:} \url{https://www.unb.ca/cic/datasets/ids-2017.html}
or via Kaggle: \texttt{kaggle datasets download -d cicdataset/cicids2017}

\section{Bot-IoT}

Koroniotis et al.\ \cite{koroniotis2019} provide IoT botnet traffic
with 5 million records across normal and four attack types (DDoS, DoS,
reconnaissance, theft).

\textbf{Download:} \url{https://research.unsw.edu.au/projects/bot-iot-dataset}

\section{Synthetic Dataset}

For CPU-only experiments, \caesar{} generates a 30,000-sample
CICIDS2017-matched synthetic dataset with feature distributions
calibrated to published statistics.
The generator is implemented in \texttt{caesar/dataset.py} and uses
a fixed seed (42) for reproducibility.
